% Chapter 2, Figure 4: vector representation of positive angular velocity components
% (scan: figures/ch2/fig04.png). The rocket, body axes and view are those of Figure 3 (ch2/fig03.tex, copied):
% a true orthographic view (elevation 32 deg, azimuth 45 deg), TikZ x = E, y = F (towards the nose), z = D.
% Each positive component is an arrow along its own axis from the centre of gravity, pointing the way a
% right-hand screw turned in the sense of Figure 3's rotation would advance (chapters/ch2-intro-sec1.tex:
% 209-213); lengths are arbitrary (the scan's proportions). omega_F lies on the rocket's axis, drawn over the
% body towards the nose, as in the scan.
\documentclass[11pt]{standalone}
\usepackage{tamrfig}
\begin{document}
\begin{tikzpicture}[x={(1.15cm,0.61cm)}, y={(-1.15cm,0.61cm)}, z={(0cm,1.38cm)}]
  % rocket: CG at the origin; nose tip at y = 3.7, nose base 2.8, tail 1.65 aft of the CG (as Figure 3)
  \def\r{0.085}\def\yn{3.7}\def\ynb{2.95}\def\yt{-1.65}
  \def\cr{0.75}\def\ct{0.3}\def\s{0.45}\def\sw{0.45}
  \pgfmathsetmacro\yle{\yt+\cr}\pgfmathsetmacro\ytle{\yle-\sw}\pgfmathsetmacro\ytte{\ytle-\ct}
  \def\phis{48.5} % silhouette angle of the body in this view
  \coordinate (s1t) at ({\r*cos(\phis)},\yt,{\r*sin(\phis)});
  \coordinate (s1n) at ({\r*cos(\phis)},\ynb,{\r*sin(\phis)});
  \coordinate (s2t) at ({-\r*cos(\phis)},\yt,{-\r*sin(\phis)});
  \coordinate (s2n) at ({-\r*cos(\phis)},\ynb,{-\r*sin(\phis)});
  \coordinate (tip) at (0,\yn,0);
  % centre line aft of the tail
  \draw[centerline] (0,\yt,0) -- (0,\yt-1.5,0);
  % fins behind the body: +E and -D
  \draw[outline, fill=white] (\r,\yle,0) -- (\r+\s,\ytle,0) -- (\r+\s,\ytte,0) -- (\r,\yt,0) -- cycle;
  \draw[outline, fill=white] (0,\yle,-\r) -- (0,\ytle,-\r-\s) -- (0,\ytte,-\r-\s) -- (0,\yt,-\r) -- cycle;
  % body and nose
  \def\nose{(s1t) -- (s1n) .. controls ($(s1n)+(0,0.4,0)$) .. (tip) .. controls ($(s2n)+(0,0.4,0)$) .. (s2n) -- (s2t)}
  \fill[white] \nose -- cycle;
  \draw[outline] \nose;
  \draw[outline, canvas is xz plane at y=\ynb] (\phis:\r) arc[start angle=\phis, end angle=\phis+180, radius=\r];
  \draw[outline, fill=white, canvas is xz plane at y=\yt] (0,0) circle[radius=\r];
  % fins in front: +D and -E
  \draw[outline, fill=white] (0,\yle,\r) -- (0,\ytle,\r+\s) -- (0,\ytte,\r+\s) -- (0,\yt,\r) -- cycle;
  \draw[outline, fill=white] (-\r,\yle,0) -- (-\r-\s,\ytle,0) -- (-\r-\s,\ytte,0) -- (-\r,\yt,0) -- cycle;
  % body axes from the centre of gravity
  \draw[thin vec] (0,0,0) -- (0,0,2.3) node[above] {D};
  \draw[thin vec] (0,0,0) -- (2.5,0,0) node[right] {E};
  \draw[thin line] (tip) -- (0,\yn+0.9,0);
  \draw[thin vec] (0,\yn+0.9,0) -- (0,\yn+1.5,0) node[above left=-1pt] {F};
  % the positive angular velocity components as vectors along their axes
  \draw[vec] (0,0,0) -- (0,0,1.2) node[right=2pt, pos=0.9] {$\omega_D$};
  \draw[vec] (0,0,0) -- (1.3,0,0) node[below right=0pt, pos=0.92] {$\omega_E$};
  \draw[vec] (0,0,0) -- (0,1.55,0) node[above=4pt, pos=0.88] {$\omega_F$};
  \node[cg mark] at (0,0,0) {};
\end{tikzpicture}
\end{document}
